\documentclass[11pt,letterpaper]{article}
\pdfoutput=1
\usepackage{jheppub}
\usepackage{mathtools,amssymb,amsthm}
\newtheorem{theorem}{Theorem}[section]
\hypersetup{pdftitle={The topological string on balanced local curves from two circle actions on a Calabi-Yau fivefold},
 pdfauthor={Bernd Johannes Wuebben},
 pdfsubject={Stable-pair indices of local Calabi-Yau fivefolds},
 pdfkeywords={topological strings, M-theory, stable pairs, local curves}}
\title{The topological string on balanced local curves\\
from two circle actions on a Calabi--Yau fivefold}
\author{Bernd Johannes Wuebben}
\affiliation{New York, NY, USA}
\emailAdd{wuebben@gmail.com}
\date{November 25, 2026}
\makeatletter
\def\@fpheader{}
\preprint{\small\shortstack[r]{\@date\\This version: November 25, 2026}}
\makeatother
% The JHEP title already breaks before the ruled contents page.
\renewcommand{\afterTocRuleSpace}{\newpage}
\abstract{In the membrane proposal of Nekrasov and Okounkov, the
variable that counts the Euler characteristic of a stable pair on a
threefold is a rotation fugacity of an ambient Calabi--Yau fivefold.
We prove the predicted exchange of these roles for two circle actions
over a smooth curve of positive genus. One fixes a local Calabi--Yau
threefold with two normal lines of equal degree. The other fixes the
product of the curve with a plane. With the orientations determined
by the ordered normal lines, their stable-pair indices agree in every
positive curve degree as integral classes on the common space of cycles,
throughout the Picard variety of splittings. On the product threefold
a deformation of the matrix equations makes the virtual contribution
vanish above the minimum Euler characteristic. At the minimum, the
rotation weights of plane Hilbert schemes reproduce the known local-curve
partition function. For a doubled genus-two curve the connected index
is supported where the two transverse points coincide. A first variation
of the rotation weights detects the same nonzero multiple of the Picard
theta class in both descriptions, although the corresponding coefficient
vanishes on every individual target.}
\keywords{Topological Strings, M-Theory, Differential and Algebraic Geometry}
\begin{document}
\maketitle
\section{Introduction}\label{sec:intro}

Fix the total wrapping charge of membranes on a curve. What accounts
for the dependence of the topological-string partition function on its
coupling at that charge? In the stable-pair calculation on a local
Calabi--Yau threefold, this dependence is encoded by a variable $q$
that grades the Euler characteristic of the sheaf: a pair of Euler
characteristic $n$ has weight $(-q)^n$. In the fivefold description
proposed by Nekrasov and Okounkov \cite{no}, the same $q$ is the
fugacity for a circle rotating the two complex directions normal to
the threefold. Selecting a different circle in the fivefold can change
the threefold on which one counts. Can the rotation weights of the new
geometry recover the Euler-characteristic series of the old one?

We answer this question for two explicit circles over a smooth curve
of positive genus. One fixes a local Calabi--Yau threefold in which the
curve has two normal line bundles of equal degree. The other fixes
the product of the curve with a plane. On the product threefold, $q$
rotates the plane, and a different variable counts the stable pairs.
We show why this second counting variable disappears and how the
weights of the plane rotation reproduce the entire $q$-dependence in
each curve degree. The two calculations agree already on the space of
curves in the fivefold, as the normal line bundle varies. This is a
case of the independence of the external circle formulated by Holmes
and Schuler \cite[Conjecture~1.7 and Examples~1.13--1.14]{holmesschuler}.

The distinction between the two variables of the topological string
will be important. We write $Q^m$ for curve degree $m$ and $q=e^{iu}$
for the exponential of the string coupling $u$. The wrapping variable $Q$ is common to both descriptions. In the Gopakumar--Vafa formula, the dependence on $u$
at fixed total degree comes from spin traces and repeated windings of
wrapped particles \cite{gv,dw}. Here we derive that dependence by
comparing two geometric indices. The doubled curve in
section~\ref{sec:collision} will let us ask how much this calculation
can tell us about a particle of charge two.

\subsection{Two threefolds in one fivefold}
Fix a smooth projective complex curve $C$ of genus $g\geq1$, put $b=g-1$,
and let $B=\operatorname{Pic}^{b}(C)$. Over $C\times B$ choose the
Poincar\'e line $\mathcal L$, trivialized along a chosen point of $C$,
and set $\mathcal M=K_C\otimes\mathcal L^{-1}$. For a line bundle $U$
and a torus character $w$, write $Uw$ for the bundle with fiber weight $w$.
Our family of fivefolds and its two fixed threefolds are
\begin{align}
 Z&=\operatorname{Tot}_{C\times B}
       (\mathcal Lt\oplus\mathcal Mt^{-1}\oplus\mathcal Oq\oplus\mathcal Oq^{-1}),
       \nonumber\\
 X=Z^{\mathbb C_q^\times}
   &=\operatorname{Tot}_{C\times B}(\mathcal L\oplus\mathcal M),
 &X'=Z^{\mathbb C_t^\times}
   &= C\times B\times S,\qquad S=\mathbb A^2 .
 \label{eq:fixed}
\end{align}
The dimensions are relative to $B$. Here a circle is a fiber rotation,
written algebraically as its complexification $\mathbb C^\times$. On $X$
the internal circle rotates the two line bundles with weights $t,t^{-1}$.
On $X'$ it rotates the plane with weights $q,q^{-1}$. Both external
circles preserve the fivefold volume form. The canonical line of $X'$ is
the pullback of $K_C$; its two normal lines in $Z$ supply the compensating
determinant.

The first fixed threefold is the balanced antidiagonal local curve used
to construct universal curve-counting classes supported on the multiples
of its zero section \cite[\S2]{wuebben-antidiagonal}. As the multiplicity
and genus vary, including genus zero, these classes freely generate the
part of Pardon's integral curve-counting algebra in first
Chern number and virtual dimension zero
\cite[Theorem~1.1]{wuebben-integral}. Embedding this model in $Z$ supplies
a second fixed threefold and a common space of cycles on which to
compare the two stable-pair indices.

A stable pair consists of a pure one-dimensional sheaf $\mathcal F$ with
proper support and a section $\mathcal O\to\mathcal F$ whose cokernel is
zero-dimensional \cite{pt}. Its curve degree $m$ is the degree of its
fundamental cycle over $C$, and $n=\chi(\mathcal F)$.
We use the modified virtual structure sheaf of \cite{no}.
The virtual structure sheaf counts with the deformation and obstruction
spaces of the pair. Its modification uses a square root of the virtual
canonical line and the determinant of cohomology from the two ordered
external lines. Exchanging the circles changes those lines as well as the
fixed threefold. We specify both roots in sections~\ref{sec:original}
and~\ref{sec:rotated}. Throughout the comparison we work relative to the
fixed curve, without integrating over $B$.

\subsection{The common index}

Curve degree one illustrates the exchange. On the product threefold a
copy of $C$ can sit at any point of the plane. At the minimum Euler
characteristic $1-g$, these translates are all the pairs. Their
deformations are the two plane translations, and their obstructions
are $g$ copies of the same two directions, one for each element of a
basis of $H^1(C,\mathcal O_C)$. The only fixed translate sits at the
origin. Cancelling the deformation factor against one obstruction
factor leaves the rotation character
\[
             Z_1(q)=\big[(1-q)(1-q^{-1})\big]^{g-1}.
\]
For a genus-two curve this is $2-q-q^{-1}$. In the original threefold
the three powers of $q$ belong to three Euler characteristics. In the
plane they all come from rotations at a single fixed pair. The
all-degree theorem extends this observation to nonreduced curves and
keeps track of their support.

Let $W_m$ be the relative Chow space of compact effective one-cycles of
degree $m$ in $Z/B$. Stable pairs on either fixed threefold determine
points of this space by taking their fundamental cycles in $Z$.
Write $s_m:B\to W_m$ for the section $m[C]$. The group
$G_0^{q,t}(W_m)$ records equivariant coherent sheaves modulo exact
sequences, with their integral multiplicities and torus characters.
Let $\mathcal I_m^{(q)}$ and $\mathcal I_m^{(t)}$ be the series of the
proper direct images of the modified sheaves in the two presentations.
Our counting convention gives weight $(-q)^n$ in the first. In the
second the weight is $(-t)^{n+mb}$, including the shift for the ordered
external line of degree $b$. The theorem in particular makes both sums
finite in this group.

\begin{theorem}\label{thm:exchange}
For the family \eqref{eq:fixed}, with the ordered determinant orientations
above, the two stable-pair indices have the same class in
$G_0^{q,t}(W_m)$ for each $m\geq 1$. More precisely,
\begin{equation}
 \mathcal I_m^{(q)}=\mathcal I_m^{(t)}
 =(s_m)_*\bigl(Z_m(q)[\mathcal O_B]\bigr),\qquad
 Z_m(q)=\sum_{\lambda\vdash m}
      \prod_{\square\in\lambda}
        \big[(1-q^{h(\square)})(1-q^{-h(\square)})\big]^{g-1}.
 \label{eq:hook}
\end{equation}
Here $h(\square)$ is one plus the number of boxes to the right and below
$\square$ in its row and column. The equality holds over the full Picard
variety, including the locus where the number of sections of the normal
bundle jumps.
\end{theorem}

On a fixed target the character $Z_m(q)$ is known. Bryan and
Pandharipande obtained the local Gromov--Witten partition function
\cite{bryanpandharipande}, and Monavari's stable-pair formulas give
its $q$-character \cite{monavari,refined}. For the related comparison
with fivefold Gromov--Witten theory, Giacchetto, Pandharipande and
Schuler derive a degree-one local-curve correspondence from their
membrane conjectures \cite{gps}. The content of
Theorem~\ref{thm:exchange} is that two moduli problems give the
same integral class supported on $m[C]$. Equality of the Laurent
polynomials on individual targets would leave this undetermined:
a nonzero class on the Picard variety can have rank zero at every
point.

The two sides account for the equality in different ways. On $X$
the infinitesimal rotation of the normal lines gives a cosection of
the obstruction sheaf. It concentrates the virtual sheaf on invariant
pairs. Their fixed deformation theory, with its determinant orientation,
is constant over $B$, even where the dimension of the space of normal
sections jumps (section~\ref{sec:original}).

On $X'$ the two plane coordinates act as endomorphisms of a rank-$m$
bundle $E$ on $C$. Its degree $k=n+m(g-1)$ is precisely the exponent
of the exchanged counting variable. Write $x,y$ for the endomorphisms.
The framed matrix equation admits the deformation
\[
                      [x,y]+ij=a\,1_E ,
\]
where $i$ is the section of the pair, $j$ a covector, and $a$ a scalar.
At $a\ne0$, taking traces on invariant quotients and subspaces forces
the words in $x,y$ applied to $i$ and to $j$ to generate $E$ and
$E^\vee$, respectively. Both bundles are therefore generated by global
sections. Their degrees must both be nonnegative, so $k=0$.
For $k>0$, specialization back to the original equation makes the
virtual structure sheaf vanish on the stable-pair space itself. At
$k=0$ the matrix data are constant along $C$ and specify a length-$m$
subscheme of the plane. The tangent weights of its Hilbert scheme
supply the hook lengths in \eqref{eq:hook}
(section~\ref{sec:rotated}). This accounts for the loss of the
$t$-grading and the survival of the $q$-dependence in the same calculation.

\subsection{What the doubled curve detects}
In genus two, degree two is the first place where two independent
degree-one objects contribute at the same total charge as a connected
double. Over two distinct points of the plane the index sheaf factors.
This prescribes the graded symmetric square to subtract for the
independent pair. The remaining class is supported where the two
points coincide. There the rotation has two fixed length-two schemes,
one for each coordinate direction in which a point can thicken. Each
contributes $+1$ at $q^{-3}$ and $-2$ at $q^{-2}$. The graded symmetric
square contributes zero at both powers, giving the connected
coefficients $2$ and $-4$ (section~\ref{sec:collision}).
To interpret the latter as the spin trace of a charge-two particle,
one would also have to establish a normalizable state in the wrapped
theory with the continuum of two separated branes open and derive its
spin.

Varying the product $\kappa$ of the normal weights on $X$ detects
the principal polarization of the Picard variety. Give the four lines
in \eqref{eq:fixed} the weights $t,\kappa/t,q,\kappa^{-1}/q$, so their
product remains one. Write $\delta=1-\kappa$.

\begin{theorem}[First variation in genus two and curve degree two]\label{thm:picard-variation}
Let $C$ have genus two, let $B=\operatorname{Pic}^1(C)$, and put
$A=R\pi_*\mathcal L$ for the normalized Poincar\'e line on
$C\times B$, with $\pi:C\times B\to B$. With the ordered orientations
of Theorem~\ref{thm:exchange}, the original and rotated family indices
have the same coefficient of $\delta q^{-3}t^k$ for every $k\geq1$.
Its value in $K^0(B)$ is
\[
       -2\bigl(1+2\,\mathbf1_{2\mid k}\bigr)\psi^k(A),
\]
where $\psi^k$ is the $k$th Adams operation. On $X$ one takes the
$q^{-3}$ coefficient and expands its first variation at $t=0$.
On $X'$ one takes the $t^k$ coefficient and expands its first
variation at $q=0$.
\end{theorem}

At $k=2$ the Chern character of this class is $12\theta$, where
$\theta$ is the principal polarization of $B$. On $X$ it comes from
two ribbons, one in each normal line. On $X'$ the two types of divisor
contribute $4\theta$ and $8\theta$. Restriction to any individual
target gives zero because $\chi(C,L)=0$ for every $L\in B$.
Section~\ref{sec:variation} gives both calculations at the lowest
original grading $q^{-3}$.

\section{The index on the split threefold}\label{sec:original}

Take \(q\) as the external circle. Its fixed threefold over
\(L\in B=\operatorname{Pic}^{b}(C)\) is the total space of
\(L\oplus K_C L^{-1}\), where \(b=g-1\). The internal circle
\(T=\mathbb C^\times_t\) rotates these two lines with opposite weights.
We shall compute the class of the index on the space of cycles.
This requires keeping the dependence on \(L\). For example, a section of
\(L\) translates the curve inside its threefold, and the number of
independent translations can jump as \(L\) varies.

Write \(\mathcal L\) for the normalized Poincar\'e line on \(C\times B\),
\(\mathcal M=K_C\mathcal L^{-1}\), and
\(X=\operatorname{Tot}_{C\times B}(\mathcal L t\oplus\mathcal M t^{-1})\).
Let \(\mathcal P=P_n(X/B,m)\) parametrize stable pairs
\(\mathbb I=[\mathcal O_X\to\mathcal F]\) in the fibers of \(X\to B\).
The sheaf has proper pure one-dimensional support, generic degree \(m>0\)
over \(C\), Euler characteristic \(n\), and a section with
zero-dimensional cokernel~\cite{pt}.
Let \(h:\mathcal P\to W_m\) send a pair to its support cycle in the
fivefold, and let \(s_m:B\to W_m\) denote the cycle \(m[C]\).
The support morphism is proper. Indeed, compactifying \(X\) to the
projective bundle with infinity divisor \(D_\infty\) gives the projective
stable-pair support morphism. Its restriction over cycles disjoint
from \(D_\infty\) parametrizes pairs on \(X\), since purity forces the
whole sheaf support to lie on its cycle. The inclusion of cycles on the
fixed threefold into the fivefold cycle space is closed.

\subsection{The orientation and the support of the index}

Push the universal sheaf down to \(C\times\mathcal P\) and call the result
\(E\). The projection of the support is proper and affine, hence finite.
Purity makes every fiber of \(E\) over \(\mathcal P\) torsion-free on \(C\).
Flatness over \(\mathcal P\) then makes \(E\) a vector bundle of rank \(m\).
For the projection \(\rho:C\times\mathcal P\to\mathcal P\), set
\begin{equation}
 P=R\rho_*E-
 R\rho_*\bigl(\operatorname{End}(E)(1-\mathcal L t)\bigr).
 \label{eq:originalpolarization}
\end{equation}
This is a polarization of the relative virtual tangent class:
\begin{equation}
 T^{\rm vir}_{\mathcal P/B}=P-P^\vee,\qquad
 \operatorname{rk}P=n+m^2b .
 \label{eq:originaltangent}
\end{equation}
Indeed the two normal coordinates give a Koszul resolution for the
self-Ext term, whose class before pushforward is
\(\operatorname{End}(E)(1-\mathcal Lt)(1-\mathcal M t^{-1})\).
Duality on \(C\) pairs its two halves and pairs the two terms linear in
\(E\) in the framed deformation complex. Riemann--Roch gives the rank
in~\eqref{eq:originaltangent}. We use \((\det P)^{-1}\) as the square
root of the relative virtual canonical bundle. The two ordered external
lines are trivial as bundles in this presentation, so their determinant
ratio introduces no further line. Write
\(\widehat{\mathcal O}^{\rm vir}=(\det P)^{-1}\mathcal O^{\rm vir}\)
with the signed counting variable \(-q\)~\cite{no}.

Let \(F=\mathcal P^T\) and \(i:F\to\mathcal P\) denote its inclusion.
A fixed pair has its entire reduced support on the zero section. Indeed,
every fiber of its support over \(C\) is finite, so the connected torus
fixes each point of that fiber. The origin is the only fixed point in
each affine-plane fiber. In particular \(F\) is proper over \(B\),
by properness of the support morphism and invariance of properness
under nilpotent thickenings. Fixed supports can themselves have
nontrivial scheme structure.

To obtain an \emph{integral} lift of the virtual sheaf to \(F\), consider
the vector field induced by the normal rotation. Relative Calabi--Yau
duality identifies \(\operatorname{Ob}_{\mathcal P/B}\) with
\(\Omega_{\mathcal P/B}\), and contraction with the vector field gives
a cosection. We check that it descends to the absolute obstruction
sheaf. Write \(\xi\) for the Euler field on \(X\) and
\(D=\iota_\xi\operatorname{At}_{X/B}(\mathbb I)\) for the corresponding
degree-one endomorphism. A Picard tangent represented by
\(\alpha\in H^1(C,\mathcal O_C)\) varies the transition functions of
\(\mathcal L\) and \(\mathcal M\) inversely. Its Kodaira--Spencer class
is therefore \(\alpha\xi\), and its obstruction is the trace-free part
of \(\alpha D\), up to the common infinitesimal-action sign~\cite{ht}.
Its pairing with \(D\) vanishes. In the Serre trace, commuting \(D\)
past the even block \(\alpha D\) preserves the sign, whereas commuting
it past the odd scalar \(\alpha\) changes the sign. This calculation
also applies over nonreduced parameter schemes. For the proper Serre
trace one can use the projective bundle compactification. The trace-free
endomorphism complex vanishes near infinity, where \(\mathbb I\) is the
identity line. The relative canonical bundle on the compactification is
\(\mathcal O(-3D_\infty)\), and the section of
\(\mathcal O(3D_\infty)\) identifies the twisted trace-free complex
with the untwisted one in the derived category. Multiplication by this
section has degree zero and leaves the cyclic sign unchanged.

Thus the cosection kills the image of \(T_B\) in the relative obstruction
sheaf and descends through the absolute obstruction sequence. Its zero
scheme is \(F\): on a homogeneous function of weight \(j\) the orbit
derivation is multiplication by \(j\), and the nonzero weights
generate the fixed-scheme ideal in characteristic zero.
Equivariant cosection localization now gives
\begin{equation}
 \gamma\in G_0^T(F),\qquad i_*\gamma=[\mathcal O_{\mathcal P}^{\rm vir}].
 \label{eq:originalintegrallift}
\end{equation}
Here \(G_0^T\) denotes the group of equivariant coherent sheaves with
integer coefficients. We use the absolute virtual sheaf induced by the
relative theory over the smooth base and the relative orientation
fixed above~\cite{kiemli,kiemsavvas}.

\subsection{The fixed class over the Picard variety}

On \(F\), decompose \(P=P_0+P_{\rm mov}\) into its zero and nonzero
weights. The moving virtual normal class is
\(P_{\rm mov}-P_{\rm mov}^\vee\).
To compute \(\gamma\) without losing torsion in the Picard family,
choose a two-term bundle resolution
\([N^{-1}\to N^0]\) of the moving cotangent complex on \(F\).
One can choose it weight by weight since \(F\) is quasi-projective and
\(T\) acts trivially on it. For \(e(V)=\lambda_{-1}(V)\), the
self-intersection and composition identities in virtual pullback
give~\cite[proof of Theorem 3.3 and Lemma 3.6]{qu}
\begin{equation}
             e(N^0)\gamma=e(N^{-1})[\mathcal O_F^{\rm vir}].
 \label{eq:originalintegralcancel}
\end{equation}
Indeed virtual self-intersection applied to \(i_*\gamma\) gives the
left side. Applied to the virtual sheaf on \(\mathcal P\), the same
pullback adjoins \(N^{-1}[1]\) to the fixed obstruction theory and
gives the right side. Both operations precede division by torus
characters: the factors in~\eqref{eq:originalintegralcancel} are the
Koszul classes for zero sections.

On each component a moving bundle has Euler polynomial in \(t\) with
invertible extremal coefficients, up to sign and a line bundle.
Hence multiplication by it on \(G_0(F)[t^{\pm1}]\) is injective
even in the presence of additive torsion. The identity
\(1-a=-a(1-a^{-1})\) applied to the lines of a splitting of the
moving polarization gives
\[
 \frac{e(N^{-1})}{e(N^0)}
     =(-1)^{\operatorname{rk}P_{\rm mov}}\det P_{\rm mov}.
\]
Clearing the Euler factors and using injectivity in
\eqref{eq:originalintegralcancel} identifies the oriented lift
integrally as
\begin{equation}
 (\det P)^{-1}\gamma=
 (-1)^{\operatorname{rk}P_{\rm mov}}
       (\det P_0)^{-1}[\mathcal O_F^{\rm vir}].
 \label{eq:originalfixedlift}
\end{equation}
In particular its internal rotation weight is zero.

It remains to determine how this class varies with \(L\).
Fix \(L_0\in B\) and put \(A=\mathcal L\otimes L_0^{-1}\).
An equivariant stable pair on the fiber at \(L_0\) gives a pair on the
varying threefold by tensoring each weight-\(j\) summand \(E_j\) with
\(A^j\). Both normal actions have precisely the required changes of
weight and the section has weight zero. The inverse construction uses
\(A^{-1}\). The degree of \(A\) on each curve is zero, so this preserves
\(n\). Stable pairs have trivial automorphisms and a fixed family has a
unique compatible linearization. The construction commutes with
arbitrary base change and gives \(F\simeq F(L_0)\times B\).

This identification includes the fixed obstruction theory. In local
frames of \(A\), the weight-zero relative derived endomorphisms, the
trace and the maps induced by the section glue with total power
\(A^0\). The same identification respects their obstruction morphism:
the relative Atiyah class in the pair-parameter directions is taken
with \(A\) held fixed on \(C\times B\).
Thus the relative fixed virtual sheaf is pulled back from \(F(L_0)\).
The terms of \(P_0\) also have total power \(A^0\), and the moving
rank in~\eqref{eq:originalfixedlift} is constant on each product
component. Proper pushforward of the oriented fixed class therefore
gives a numerical coefficient \(c_{m,n}\) times \([\mathcal O_B]\).
On \(F_{\rm red}\) the support map factors through \(s_m\). Filtering
coherent sheaves by the nilradical permits the same pushforward
calculation even if the map on \(F\) has infinitesimal components
normal to \(s_m(B)\). Consequently
\begin{equation}
  h_*[\widehat{\mathcal O}_{\mathcal P}^{\rm vir}]
     =(s_m)_*\bigl(c_{m,n}[\mathcal O_B]\bigr)
     \quad\text{in }G_0^T(W_m).
 \label{eq:originalsupport}
\end{equation}

It suffices now to evaluate \(c_{m,n}\) on one fiber. The antidiagonal
local-curve stable-pair formula at levels \((b,b)\)
gives~\cite{monavari,refined}
\begin{equation}
 \sum_n c_{m,n}(-q)^n
   =\sum_{\lambda\vdash m}\ \prod_{u\in\lambda}
     \bigl((1-q^{h(u)})(1-q^{-h(u)})\bigr)^b ,
 \label{eq:originalhook}
\end{equation}
where \(h(u)\) is one plus the number of boxes to the right and below
\(u\) in its row and column. The sign follows from
\((1-q^h)(1-q^{-h})=-(q^{h/2}-q^{-h/2})^2\). In degree one this
gives \(((1-q)(1-q^{-1}))^b\). For \(g\geq1\) every summand is an
integral Laurent polynomial. Equation~\eqref{eq:originalsupport}
identifies the family class to which this numerical formula belongs.

\section{The rotated threefold}\label{sec:rotated}

Choose the circle with character \(t\) as the external circle. Its fixed
threefold is \(X'=C\times S\), where \(S=\mathbb A^2\) carries the two
weights \(q,q^{-1}\). A curve can now move off the origin of \(S\) by a
constant translation. We compute on this noncompact moduli space before
pushing its index to cycles.

\subsection{Deforming the framed matrix equation}

Let \(\mathcal O_{X'}\to F\) be a stable pair of curve degree \(m>0\) and
Euler characteristic \(n\). Its proper support projects finitely to \(C\).
Indeed, the projection is affine and its restriction to the support is
proper. Purity then makes \(E=p_*F\) a vector bundle of rank \(m\),
where \(p:X'\to C\). Riemann--Roch gives
\begin{equation}
 k=\deg E=n+m(g-1).
 \label{eq:rotated-excess}
\end{equation}
The two coordinates of \(S\) act on \(E\), and the section of the pair
gives a distinguished vector \(i:\mathcal O_C\to E\). Over the function
field of \(C\), the vector generates \(E\) under the coordinate actions.

The framed ADHM description of the stable-pair theory adds a covector
\(j:E\to\mathcal O_C\) and writes the relation as a matrix equation
\cite[Theorem~1.9 and \S7]{diaconescu}. Write \(x,y\) for the two
endomorphisms. We deform the equation, with its scalar part included, to
\begin{equation}
 [x,y]+ij=a\,1_E,\qquad a\in\mathbb A^1 .
 \label{eq:rotated-adhm}
\end{equation}
Framing fixes the copy of \(\mathcal O_C\), and we divide by changes of
basis of \(E\). The stable open is specified by generic generation by
the words in \(x,y\) applied to \(i\). At \(a=0\) this is
Diaconescu's presentation of the stable-pair moduli. Under the opposite
weights on \(x,y\), the parameter \(a\) is invariant.

At \(a\ne0\), fix any point of \(C\). The span of the words applied to
\(i\) is invariant under \(x,y\). On its quotient the commutator would be
\(a\) times the identity. Taking the trace forces this quotient to vanish.
Dually, the common kernel of the covectors \(j\,w(x,y)\), for all words
\(w\), is invariant and contained in \(\ker j\). The same trace argument
on this subspace makes it zero. Words of length at most \(m-1\) suffice
in each argument, since an increasing span which stops growing is already
invariant. These words are global sections on \(C\). Consequently \(E\)
and \(E^\vee\) are both globally generated and \(\deg E=0\). At \(a=0\)
there are still \(m\) words generically independent on \(i\). Their
determinant is a nonzero section of \(\det E\), so \(k\ge0\).

For \(k>0\) this calculation implies the integral identity
\begin{equation}
 [\mathcal O^{\mathrm{vir}}_{P_n(X',m)}]=0
 \qquad\text{in }G_0^{\mathbb C^\times_q}(P_n(X',m)).
 \label{eq:rotated-vanishing}
\end{equation}
We give the specialization argument, since the equivariant weight of
\(a\) matters here. Use the derived zero locus of
\eqref{eq:rotated-adhm} modulo changes of frame and take its stable space
of sections over the fixed curve. Relative to the \(a\)-line its tangent
complex is the derived push from \(C\) of the complex in degrees
\(-1,0,1\)
\begin{equation}
 \operatorname{End}E
 \longrightarrow (\operatorname{End}E)q\oplus
  (\operatorname{End}E)q^{-1}\oplus E\oplus E^\vee
 \longrightarrow \operatorname{End}E .
 \label{eq:rotated-deformation}
\end{equation}
The first arrow is the infinitesimal change of basis and the second is
the differential of the moment map. Generic cyclicity makes the first
injective over the function field. The trace symplectic pairing on the
framed matrix space identifies its transpose, up to sign, with the
second arrow, which is therefore generically surjective. On \(C\) the
first kernel sheaf vanishes and the last cokernel has finite support.
Pushing to the parameter space gives tangent amplitude \([0,1]\).
This is the fixed-domain quasimap obstruction theory
\cite[\S4.3]{okounkov}. Stabilizers are trivial by generic cyclicity.
Boundedness at fixed \(k\) follows by dualizing \(m\) generically
independent word sections. The resulting subbundle away from finitely
many points, \(E^\vee\subset\mathcal O_C^m\), has quotient of length
\(k\), hence belongs to a finite-type Quot scheme.

For clarity, the comparison on the central fiber can be made at the
level of virtual structure sheaves. If \(\pi\) forgets the \(C\)-factor
of the universal family, the stable-pair virtual tangent has class
\begin{equation}
 R\pi_*\!\left(E+E^\vee-
       (\operatorname{End}E)(1-q)(1-q^{-1})\right).
 \label{eq:rotated-tangent}
\end{equation}
The Koszul resolution for the two plane coordinates gives the quadratic
term in \(E\) in the trace-free pair complex. The two linear terms come
from the section and its dual, with the product of the two plane weights
equal to one. This agrees with the alternating class of
\eqref{eq:rotated-deformation} on the same moduli scheme. The
\(K\)-theoretic Fulton class formula expresses the virtual structure
sheaf of a quasi-projective scheme in terms of its intrinsic Fulton
class and its virtual cotangent class
\cite[Theorem~4.2 and Corollary~4.5]{thomas}. Its cone filtration and
exterior powers preserve the \(q\)-linearizations, with a separate
auxiliary grading for the filtration. Thus the comparison holds in
integral equivariant coherent \(K\)-theory. See also the local-curve
application in \cite[\S9]{monavari}.

The virtual structure sheaf on the central fiber is the refined pullback
of the relative virtual family along \(a=0\), by derived base change
or virtual specialization \cite{qu}. For \(k>0\) the underlying space
of the entire family lies over zero. Filter a coherent sheaf there by
powers of \(a\). Every successive quotient is a sheaf on the central
fiber, and its refined pullback has two equal Tor terms in adjacent
degrees, since \(a\) has weight zero. The pullback thus kills every
class of the total family. This proves \eqref{eq:rotated-vanishing},
even when the parameter acts nilpotently on its ordinary structure
sheaf.

\subsection{The plane Hilbert scheme and its orientation}

At \(k=0\) the determinant of the \(m\) independent word sections has
degree zero and no zero. These sections trivialize \(E\) on \(C\).
The coordinate operators and the cyclic vector of the stable pair are
then constant, and give a cyclic length-\(m\) quotient of the coordinate
ring of \(S\). This identification works in families. Locally on the
parameter scheme choose words that form a basis on one fiber over a
point of \(C\). At the chosen parameter they form a basis on the entire
curve. Properness of \(C\) allows one to shrink the parameter scheme
until they form a basis everywhere. Since \(C\) is projective and
connected, the matrix entries then descend to the parameter scheme,
also over nonreduced bases. It follows that
\begin{equation}
 P_{-mb}(X',m)\simeq H_m=\operatorname{Hilb}^m(S),
 \qquad
 \xi\longmapsto [\mathcal O_{X'}\longrightarrow
                          \mathcal O_{C\times\xi}] .
 \label{eq:rotated-hilbert}
\end{equation}

For these constant cyclic data the framed complex on \(C\) is pulled
back from the tangent bundle \(T H_m\). The virtual tangent is
\(R\Gamma(C,\mathcal O_C)\otimes T H_m\). Its degree-zero cohomology is
the ordinary tangent bundle of the smooth Hilbert scheme. Its
obstruction bundle is \(H^1(C,\mathcal O_C)\otimes T H_m\). The
symplectic form on \(H_m\) is invariant under the antidiagonal circle
and trivializes the tangent determinant. We use the resulting relative
orientation, suppressing fixed one-dimensional vector spaces from \(C\),
which have trivial circle character.

There is also a determinant contribution from the two external lines
\cite{no}. Let \(V\) be the rank-\(m\) tautological bundle on \(H_m\),
let \(\pi_B:C\times B\to B\), and put \(A=R\pi_{B*}\mathcal L\).
Relative duality gives \(R\pi_{B*}\mathcal M=A^\vee[-1]\).
The unweighted ratio of the two determinants on \(H_m\times B\) is
\begin{equation}
 D=\frac{\det(V\otimes A)}{\det(V\otimes A^\vee[-1])}.
 \label{eq:rotated-external-determinant}
\end{equation}
Here both complexes on \(B\) have virtual rank zero and the same
determinant. Thus \(D\) is trivial and we take its trivial square root.
This argument uses perfect complexes throughout, including where the
individual spaces of sections of \(\mathcal L\) jump. The external
prefactor is \((-t)^{n+mb}\), hence is one on \(H_m\). The modified
class on the minimum component is therefore the pullback from \(H_m\)
of
\begin{equation}
 \mathcal F_m=\lambda_{-1}(T^\vee H_m)^g,\qquad
 \lambda_{-1}(U)=\sum_i(-1)^i[\wedge^i U].
 \label{eq:rotated-exterior}
\end{equation}

The effect of balance can be tested on an actual family of coincident
points. With \(B=\operatorname{Pic}^d C\) set \(\delta=d-(g-1)\).
The two ranks in \eqref{eq:rotated-external-determinant} become
\(\delta,-\delta\), whereas their determinants still agree. Thus
\(D=(\det V)^{2\delta}\) with oriented root \((\det V)^\delta\).
The punctual Hilbert square is a \(\mathbb P^1\) of one-dimensional
quotients of the cotangent plane at the origin. On this line the
length-two algebra has the constants and the universal quotient
\(\mathcal O(1)\) as its two graded pieces. The degree of the root on
this collision family is \(\delta\). The choice \(d=g-1\) is therefore
detected by the disappearance of the twist on the Hilbert scheme
itself.

\subsection{Comparison on the space of cycles}

The \(q\)-fixed ideals of \(H_m\) are the monomial ideals
\(\xi_\lambda\) for partitions \(\lambda\) of \(m\). To see that the
one-parameter subgroup has no further fixed components, first observe
that every fixed zero-cycle of \(S\) is concentrated at the origin.
The punctual Hilbert scheme is projective, so every connected component
of the fixed locus contains a point fixed by the full coordinate torus.
Those points are monomial ideals. At \(\xi_\lambda\) the tangent weights
restrict to
\begin{equation}
 \big\{q^{h(s)},q^{-h(s)}:s\in\lambda\big\},
 \label{eq:rotated-hook-weights}
\end{equation}
by the arm and leg formula \cite[Lemma~5.4.5]{haiman}. They are all nontrivial, so
each such point is isolated and reduced in the fixed locus.

The infinitesimal \(q\)-rotation is an equivariant section of \(T H_m\).
Its linearization at a fixed ideal is invertible by
\eqref{eq:rotated-hook-weights}. Its zeros are therefore simple and its
Koszul complex resolves their structure sheaf. Since \(g\ge1\), one
exterior Euler factor in \eqref{eq:rotated-exterior} resolves the fixed
points. The other \(g-1\) factors restrict to them. This proves the
identity
\begin{equation}
 \mathcal F_m
   =\sum_{|\lambda|=m}[\mathcal O_{\xi_\lambda}]
        \prod_{s\in\lambda}
            \big((1-q^{h(s)})(1-q^{-h(s)})\big)^{g-1}
 \quad\text{in }G_0^{\mathbb C^\times_q}(H_m).
 \label{eq:rotated-integral-concentration}
\end{equation}
In particular no character denominator is needed for the transverse
plane, even for an elliptic curve, when each point contributes one.

Let \(s_m:B\to W_m\) be the section of the common fivefold cycle space
which assigns the zero-section cycle \(m[C]\). Hilbert--Chow sends each
\(\xi_\lambda\) to the origin of \(\operatorname{Sym}^m S\) with
multiplicity \(m\), so its associated curve has cycle \(m[C]\).
These pushforwards are proper on the displayed representatives.
Equivalently one can embed both threefold cycle spaces in \(W_m\).
Compactification by compatible projective bundles makes containment
in each closed subbundle a closed condition on cycles, and one then
restricts to the open locus disjoint from infinity.

Writing \(Z_m(q)\) for \eqref{eq:hook}, the rotated pushforward is
\((s_m)_*(Z_m(q)\mathcal O_B)\) in \(G_0^{q,t}(W_m)\). In the original presentation
section~\ref{sec:original} gives the same class after summing the
Euler-characteristic grading with its external weight. This proves
the circle exchange of Theorem~\ref{thm:exchange}. One grading has
become an ordinary rotation of the plane, and its hook weights now
come from the tangent spaces to punctual ideals.

\section{The connected index at a collision}\label{sec:collision}

Set $g=2$. The plane Hilbert square parametrizes both pairs of distinct
points and length-two schemes supported at a single point. To extract the
connected contribution we compare its index with the symmetric square
of the degree-one index, on the space of cycles itself. All constructions below can be pulled back
over $B$; we suppress its trivial structure-sheaf factor.

Let $h:H=\operatorname{Hilb}^2(S)\to W=\operatorname{Sym}^2(S)$ be
the Hilbert--Chow morphism, $p:S^2\to W$ the finite quotient and
$i:S\to W$ the map $x\mapsto 2[x]$. Put
\begin{equation}
 K_1=\lambda_{-1}(T_S^\vee)^2,\qquad
 K_2=\lambda_{-1}(T_H^\vee)^2 .
\end{equation}
Here $\lambda_{-1}(U)=\sum_j(-1)^j[\bigwedge^j U]$. Its even and odd
summands give the signs in a permutation of two such factors. The
connected class is
\begin{equation}
 D_+=h_*K_2-\bigl(p_*(K_1\boxtimes K_1)\bigr)^{\mathfrak S_2}
                    \quad\in G_0^q(W).
 \label{eq:connected-class}
\end{equation}
The invariant summand is exact over $\mathbb C$. Over two distinct points,
$h$ is an isomorphism, and on the ordered cover the tangent bundle is the
direct sum of the tangent bundles at the two points. Exchanging them acts
on the exterior algebra with the graded permutation sign. Thus the two
terms of \eqref{eq:connected-class} restrict to the same class there.
The localization sequence for coherent sheaves makes $D_+$ the pushforward
of an integral class on $i(S)$.

Write $y=q+q^{-1}$, $e=(1-q)(1-q^{-1})=2-y$, and
$e_2=(1-q^2)(1-q^{-2})=4-y^2$. The degree-one index is $e$.
At the two ideals $(z_1^2,z_2)$ and $(z_1,z_2^2)$, in plane coordinates
$z_1,z_2$, the four tangent weights on $H$ are
$q,q^{-1},q^2,q^{-2}$. Equivariant localization divides out one
cotangent Euler factor and leaves $ee_2$ at each ideal.
Their sum is $2e^2(4-e)$, the disconnected genus-two double already
obtained from the antidiagonal Gromov--Witten evaluation
\cite[\S6]{wuebben-antidiagonal}.
For the independent pair the identity permutation has trace $e^2$;
the transposition has trace $\psi^2(e)=e_2$, where the Adams operation
$\psi^2$ replaces $q$ by $q^2$. Taking invariants averages these traces.
The connected index is therefore
\begin{align}
 F_2(q)=\chi_q(D_+)
 &=2ee_2-\tfrac12(e^2+e_2)
   =2(y-2)(y^2-3) \nonumber\\
 &=2q^{-3}-4q^{-2}+4-4q^2+2q^3 .
 \label{eq:double-character}
\end{align}
The lift to the collision diagonal is determined as well.
Equivariant homotopy invariance identifies $G_0^q(S)$ with the Laurent
character ring times $[\mathcal O_S]$. The localized Euler characteristic
of $A(q)[\mathcal O_S]$ is $A(q)/e$. Since this map is injective, the
class constructed in \eqref{eq:connected-class} is
\begin{equation}
                  D_+=i_*\big(e F_2(q)[\mathcal O_S]\big) .
\end{equation}
Each punctual ideal contributes $+1$ at $q^{-3}$ and $-2$ at
$q^{-2}$. The graded symmetric square is $(e^2+e_2)/2=4-2y$.
At $q^{-2}$ the identity and the transposition have coefficients $+1$
and $-1$, respectively, so their average vanishes. Thus the two
punctual ideals account for both leading coefficients of the connected
index. The subtraction changes the constant and the $q^{\pm1}$ terms.

For completeness one can make the connected subtraction in all degrees
at the level of characters. Introduce a curve variable $Q$ and put
$Z=1+\sum_{m\ge1}Q^m Z_m(q)$. Define integer Laurent polynomials
$F_m=\sum_j a_{m,j}q^j$ successively by
\begin{equation}
 Z=\prod_{m\ge1}\prod_{j\in\mathbb Z}(1-Q^m q^j)^{-a_{m,j}} .
 \label{eq:euler-product}
\end{equation}
After dividing by the previously determined factors, the coefficient
of $Q^m$ fixes $F_m$ integrally. Each $Z_m$ is invariant under inversion
of $q$, so uniqueness gives $F_m\in\mathbb Z[2-q-q^{-1}]$.
For $g>1$ each $Z_m$ is divisible by $e^{g-1}$, since every hook Euler
factor is divisible by $e$. The same divisibility holds for $F_m$ by
induction. Indeed the nonconstant contributions from lower degrees in
the symmetric-algebra expansion contain an Adams image of a previous
$F_d$, and $2-q^r-q^{-r}$ is divisible by $e$. This proves divisibility
over $\mathbb Q$; division by the monic polynomial $e^{g-1}$ in
$\mathbb Z[e]$ proves it integrally. For $g=1$ the hook sum gives the
ordinary partition numbers and the Euler product has $F_m=1$.

Writing $F_m=\sum_{h\ge1} n_{m,h} e^{h-1}$ and putting $q=\exp(iu)$
gives the local Gopakumar--Vafa form \cite{gv,bryanpandharipande},
\begin{equation}
 \log Z=\sum_{m,h,r\ge1} \frac{n_{m,h}}{r}
        \bigl(2\sin(ru/2)\bigr)^{2h-2} Q^{mr}.
\end{equation}
In the example \eqref{eq:double-character}, $F_2=-2e+8e^2-2e^3$.
The construction \eqref{eq:connected-class} realizes this
symmetric-algebra subtraction by an actual support calculation in
coherent $K$-theory.

\section{Varying the threefold volume character}\label{sec:variation}

Moving away from the volume-preserving torus of the threefold gives a
test of circle exchange that the numerical index of a single curve cannot
detect. Let \(C\) have genus two, \(B=\operatorname{Pic}^1(C)\), and
\(\pi:C\times B\to B\). Choose a Poincar\'e line \(\mathcal L\) trivialized
over one point of \(C\), and put
\(\mathcal M=K_C\mathcal L^{-1}\), \(A=R\pi_*\mathcal L\).
Give the four lines \(\mathcal L,\mathcal M,\mathcal O,\mathcal O\)
weights \(t,\kappa/t,q,\kappa^{-1}/q\), in that order. Their product
remains one on the fivefold. Write \(\delta=1-\kappa\) and retain the
orientation of the ordered external lines in the modified pair index.

Let \(\mathcal I_2^{\rm orig}\) and \(\mathcal I_2^{\rm rot}\) denote
the degree-two family characters with these weights in the original
and rotated presentations. We prove Theorem~\ref{thm:picard-variation}
by computing both sides of the identity
\begin{equation}
 [\delta q^{-3} t^k]\,\mathcal I_2^{\rm orig}
 =[\delta q^{-3} t^k]\,\mathcal I_2^{\rm rot}
 =-2\bigl(1+2\,\mathbf 1_{2\mid k}\bigr)\psi^k(A)
 \quad\text{in }K^0(B),
 \label{eq:circle-volume-edge}
\end{equation}
for every \(k\geq1\), where \(\psi^k\) is the Adams operation. In the original presentation
one takes the \(q^{-3}\) term, varies \(\kappa\), and expands the
result at \(t=0\). After exchange, \(t\) counts pairs. One fixes \(k\),
takes the variation, and expands at \(q=0\). In the latter expansion
there are no powers below \(q^{-3}\) at this order. We give the two
calculations because the fixed loci that produce the same class
have different dimensions.

\subsection{The two minimal ribbons}

At the lowest Euler characteristic, \(n=-3\), a torus-fixed degree-two
pair is the structure sheaf of a ribbon in either normal direction.
Indeed its generic cyclic quotient of length two has monomial basis
\(1,x\) or \(1,y\). The corresponding ribbon in a line \(F\) has
pushforward \(\mathcal O_C\oplus F^{-1}\), of Euler characteristic
\(-3\). A zero divisor of the section or of the multiplication between
the two weight lines raises the Euler characteristic. Thus each
choice gives one fixed copy of \(B\).

These two ribbons also occur as limits of smooth covers. For a
genus-two splitting as in \cite[Proposition~8.1]{wuebben-geometric},
the degree-two cycles form two one-parameter families, one in each
normal line. Their nonzero members are smooth genus-four curves.
In a fiber coordinate \(z\) on either line, the family has equation
\(z^2=a s\), where \(s\) is a section of the square of that line
with two simple zeros. At \(a=0\) the central ribbon has equation
\(z^2=0\). The index calculation below keeps its scheme structure
and its equivariant deformation weights.

Write \(F=\mathcal L,\ p=t\) on the first ribbon and
\(F=\mathcal M,\ p=\kappa/t\) on the second. The equations for the
ribbon are the other normal coordinate and the square of the first.
Tensoring the corresponding normal lines with
\(\mathcal O_C\oplus(pF)^{-1}\) and using relative Serre duality gives
the virtual tangent \(H-\kappa H^\vee\), with polarization
\[
 H=p R\pi_*F+p^2 R\pi_*F^2,\qquad \operatorname{rk}H=1 .
\]
For this polarization the symmetrized inverse Euler class, the
external determinant and the sign of \((-q)^{-3}\) give
\[
 \kappa^{-1}
 \frac{\lambda_{-1}(\kappa^{-1}H)}{\lambda_{-1}(H)} ,
 \qquad
 \lambda_{-1}(V)=\sum_j(-1)^j\mathop{\bigwedge}\nolimits^j V .
\]
The ratio at \(\kappa=1\) is identically one, so varying \(p=\kappa/t\)
in both its arguments contributes zero. A root \(w\) of \(H\)
contributes \(-w/(1-w)\) to its first variation. The expansions for
\(w=tx\) and \(w=t^{-1}x\) have opposite signs in positive degree
and involve \(x^k\) and \(x^{-k}\), respectively. Using
\(R\pi_*\mathcal M=-A^\vee\), the resulting positive coefficients are
\begin{equation}
 -2\psi^k(A)+\mathbf1_{2\mid k}\,
       \psi^{k/2}(D_M^\vee-D_L),\qquad
 D_L=R\pi_*\mathcal L^2,\quad D_M=R\pi_*\mathcal M^2 .
 \label{eq:circle-unsimplified-edge}
\end{equation}
Here and below the second term is included only for even \(k\).

\subsection{The fixed pairs after circle exchange}

The rotated threefold is \(C\times\mathbb A^2\). Write
\(u=q,\ v=\kappa^{-1}/q,\ c=uv\), so \(1-c^{-1}=\delta\).
A fixed row is specified by effective divisors \(P,Q\) on \(C\).
The section vanishes on \(P\) in the inner weight line, and the
map to the outer line vanishes on \(Q\). The pushforward to \(C\) and
the counting degree are
\[
 E=\mathcal O(P)\oplus u^{-1}\mathcal O(P+Q),
 \qquad k=2\deg P+\deg Q .
\]
The transposed row exchanges \(u\) and \(v\). Thus the fixed loci are
products of symmetric powers of \(C\). The virtual tangent of a row
is the pushforward along \(C\) of the following character, obtained
by resolving the pair in the two plane directions:
\[
 E+c E^\vee-(1-u)(1-v)\operatorname{End}E .
\]
For a universal divisor \(D\) of length \(d\), let \(p_D\) be its
projection to \(S_d=\operatorname{Sym}^d C\). Set
\(U_D=p_{D*}\mathcal O_D\) and
\(T_D=p_{D*}\mathcal O_D(D)=T_{S_d}\). The divisor sequences give
the weight-one and weight-\(c\) terms
\(T_P+T_Q-c(U_P+U_Q)\). The unit in \(U_D\) splits by the normalized
trace for \(d>0\). Each nonempty divisor therefore supplies a factor
\(1-c^{-1}\). All remaining moving weights are nontrivial at \(c=1\)
for generic \(q\). When both divisors are nonempty the first
variation vanishes, including along their incidence locus.

It remains to divide by the one factor \(1-c^{-1}\) on a pure
component and set \(c=1\). Write \(U=U_D,\ T=T_D,\ J=\det U\), and
write \(N(F)\) for the norm to \(S_d\) of a line \(F\) restricted
to \(D\). Finite duality gives
\(\det T=J^{-1}N(K_C)^{-1}\). The residual fixed obstruction
contributes \(\lambda_{-1}((U/\mathbf1)^\vee)\).
The moving tangent and modified root on the two components are
\begin{equation}
\begin{array}{c|rrrr|l}
 &q^{-2}&q^{-1}&q&q^2&\text{modified root}\\ \hline
 P=0,\ Q=D&T-1&-T-1&U-1&-U-1&q^d J N(\mathcal L)\\
 P=D,\ Q=0&-1&T-1&-U-1&-1&q^d J^2 N(\mathcal L^2).
\end{array}
\label{eq:circle-pure-moving}
\end{equation}
For example, subtract the minimal tangent
\(-(u+v+u^2+v/u)\) from the virtual tangent above. The divisor
sequences give \(T(1+v/u-v)+Uu(1-u-v)\) for the first row and
\(T(1+u^{-1})-cU(1+u)\) for the second.

For clarity, the orientation in the last column can be computed
before restricting to a fixed component. Write
\(\mathcal D(F)=\det R\pi_*F\), with the projection now taken over
the pair family, and let \(\langle\, ,\,\rangle\) be the Deligne
pairing on the curve. The modified root, up to constant
vector-space lines, is
\[
 c^{k+1}\mathcal D(E)^{-1}
             \langle\det E,\mathcal L K_C^{-1}\rangle .
\]
Indeed the endomorphism term multiplied by \((1-u)(1-v)\) has
trivial determinant. For rank-two \(E\) and a line \(F\) one has
\(\mathcal D(EF)/\mathcal D(E)
=\langle\det E,F\rangle
 (\mathcal D(F)/\mathcal D(\mathcal O))^2\).
This identity and Serre duality
\(\mathcal D(E^\vee)=\mathcal D(EK_C)\),
\(\mathcal D(\mathcal L)=\mathcal D(\mathcal M)\)
give the stated root on squaring, with external character exponent
\((\chi(E\mathcal M)-\chi(E^\vee))/2=k+1\).
Restriction to the universal divisors gives the norm lines in
\eqref{eq:circle-pure-moving}. The counting sign \((-1)^k\) is
still to be included.

A summand \(q^a V\) of the virtual tangent with \(a>0\) and virtual
rank \(r\) has inverse Euler class with leading term
\((-1)^r q^{ar}\det V\) at \(q=0\). A negative weight has leading
term one. The two rows of \eqref{eq:circle-pure-moving} start at
\(q^{-3}\), since the exponent from the root and positive weights is
\(d+(d-1)+2(-d-1)=-3\) in the first and
\(d+(-d-1)+2(-1)=-3\) in the second. At \(q=\infty\) the extreme
exponent is \(3\) for each row, which gives the same lower bound
for their transposes. To extract the coefficient put
\[
 \Phi_d(F)=R\pi_{S_d*}\bigl(N(F)\lambda_{-1}(U/\mathbf1)\bigr),
 \qquad \pi_{S_d}:S_d\times B\longrightarrow B .
\]
Combining the leading determinants with
\(J\lambda_{-1}((U/\mathbf1)^\vee)
=(-1)^{d-1}\lambda_{-1}(U/\mathbf1)\) and the counting sign gives
\(-\Phi_d(\mathcal L)\) from each outer divisor, at \(k=d\).
The common divisor and its transpose, at \(k=2d\), give respectively
\(-\Phi_d(\mathcal L^2)\) and
\(-\Phi_d(\mathcal L^2K_C^{-1})\). The factor \(N(K_C)^{-1}\)
in the transpose comes from \(\det T\) at the opposite end of
the \(q\)-expansion.

\subsection{Adams operations and the Picard polarization}

These symmetric-product indices admit an integral evaluation,
\begin{equation}
 \Phi_d(F)=\psi^d(R\pi_*F).
 \label{eq:circle-symmetric-adams}
\end{equation}
To see it, let \(h:C^d\to S_d\) be the ordered cover, let
\(V=\mathbb C^d\) be the permutation representation, and put
\(V_0=V/\mathbf1\). Marking one point of a divisor gives
\(U=(h_*\mathcal O\otimes V)^{\mathfrak S_d}\).
Exterior powers of these covariants identify with the covariants of
the exterior powers. Locally at a divisor of multiplicities \(d_i\)
this is Solomon's exterior algebra theorem for the product
\(\prod_i\mathfrak S_{d_i}\) acting on the local coordinates of the
ordered points \cite{solomon}. Splitting off the constant vector
gives the same statement for \(U/\mathbf1\) and \(V_0\).
The norm pulls back to \(F^{\boxtimes d}\). Relative K\"unneth now
expresses the pushforward in \eqref{eq:circle-symmetric-adams}
as the alternating sum of the Schur operations corresponding to
\(\bigwedge^jV_0\), with the usual Koszul signs on the tensor powers
of \(R\pi_*F\). Their partitions are the hooks \((d-j,1^j)\).
The integral symmetric-function identity
\(p_d=\sum_{j=0}^{d-1}(-1)^j s_{(d-j,1^j)}\)
identifies the sum with the Adams operation. Together with
\(D_M^\vee=-R\pi_*(\mathcal L^2K_C^{-1})\), this gives
\eqref{eq:circle-unsimplified-edge} directly in the rotated
presentation.

There is one more simplification in even counting degree. On
\(S_2\), \(J\) is the trace-free line of the ordered double cover,
so \(J^2=\mathcal O(-i(C))\) for the doubled-point divisor
\(i:C\to S_2\). Its restriction \(i^*J=K_C\) is the conormal
of the diagonal on the ordered cover. Multiplication by a line
minus one lowers the dimension of the support filtration on the
smooth surface \(S_2\). Hence \((J-1)^3=0\) in \(K^0(S_2)\), and
the divisor sequence gives
\[
 4(J-1)=(J-1)(2+J+J^{-1})
        =(J^2-1)(1+J^{-1})=-i_*(1+K_C^{-1}).
\]
Tensor with \(N(\mathcal L)\) and push to \(B\). Since \(U/\mathbf1=J\)
and \(i^*N(\mathcal L)=\mathcal L^2\), equation
\eqref{eq:circle-symmetric-adams} gives
\(4\psi^2(A)=R\pi_*(\mathcal L^2(1+K_C^{-1}))\).
Applying \(\psi^{k/2}\) proves \eqref{eq:circle-volume-edge}
integrally.

Finally let \(\theta\) be the principal polarization, with sign
fixed by \(\pi_*(\gamma^2)/2=-\theta\) for the mixed K\"unneth
component \(\gamma\) of \(c_1(\mathcal L)\). The degree-one line has
Euler characteristic zero on \(C\), and the chosen normalization
removes the base component of its first Chern class. Riemann--Roch
therefore gives \(\operatorname{ch}(A)=-\theta\) in \(H^{\mathrm{even}}(B,\mathbb Q)\). The cohomological Chern
character of \eqref{eq:circle-volume-edge} is
\(2k(1+2\,\mathbf1_{2\mid k})\theta\). At \(k=2\) the two outer
components give \(4\theta\) and the common-divisor components give
\(8\theta\). Their sum \(12\theta\) measures the response of the
family to varying the threefold volume character. Restriction to
any one target retains only its rank, which is zero.

\section{Discussion}\label{sec:discussion}

In the family \eqref{eq:fixed}, changing the external circle exchanges
an Euler-characteristic grading with a rotation of a plane. The commutator
deformation accounts for the disappearance of positive excess in the
plane description. The geometry that remains is the Hilbert scheme of
points, and its tangent representation supplies the hook lengths in
the local topological-string partition function. The determinant on the
family of doubling directions distinguishes the balanced splitting:
its degree is zero precisely when the two original normal lines have
equal degree. The first variation in section~\ref{sec:variation}
illustrates what a family retains even when a numerical trace vanishes.

An index on a space of cycles and the Hilbert space of a charged
five-dimensional particle answer different questions. In degree two the
index calculation identifies the locus supporting the connected term
and fixes all its coefficients. A microscopic membrane calculation would
have to determine the spectrum at the threshold where two branes on $C$
can separate. To identify the coefficient at $q^{-2}$ with a particle
spin trace, one would have to establish a normalizable state in the full
unframed theory with the two-particle continuum open, and then derive
its spin from the wrapped-brane dynamics. The value $-4$ computed here
provides an independent comparison for that calculation.

\bibliographystyle{JHEP}
\bibliography{references}
\end{document}
